\par\Needspace{7\baselineskip}\bigskip
\begin{center}
\pdfbookmark[2]{Gauss an Olbers, Göttingen, 12. August 1811}{pallas-brief-18110812}
\textsc{Gauss} an \textsc{Olbers}, Göttingen, 12. August 1811.
\end{center}
\noindent\hspace*{\parindent}. . . . .\quad Ich habe Ihnen . . . .\ nicht früher geschrieben, um doch noch einiges über Ihre Pallas hinzufügen zu können. Einiges, was Sie in einer der nächsten Ausgaben unserer Gel. Anz. finden werden\HistoricalNoteEight{Band VI, S. 327.]} (ich glaube in der vom 17. Aug.), übergehe ich, um nur einiges von den allgemeinen Ergebnissen der Störungsrechnung zu erwähnen. In der ersten Rechnung habe ich die Störungen der Neigung und der Länge der Knotenlinie (als den am wenigsten mühsamen Teil) fertig berechnet und auch die ersteren an den 7 Oppositionen (für die meine früheren Arbeiten die Elemente angeben) geprüft. Anfangs schlechte Übereinstimmung, bis ich feststellte, dass 2 Gleichungen das falsche Vorzeichen hatten. Jetzt stimmen sie zu meiner Zufriedenheit überein. Mit der Prüfung des \(\AscendingNode\) werde ich auch in einigen Tagen fertig sein, dann will ich die Exzentrizität in Angriff nehmen, eine Arbeit von etwa einem Monat. Um Ihnen doch eine Vorstellung von den Ergebnissen zu geben, füge ich die 40 Störungsgleichungen für die Neigung bei. Ich habe nur beibehalten, was über eine Bogensekunde hinausging. Auch können vielleicht noch einige kleine Korrekturglieder hinzukommen, die vom 9-Fachen und 10-Fachen der Jupiterlänge abhängen\HistoricalNoteDouble{Hier folgen die auf S. 518 dieses Bandes angegebenen Glieder von \(\delta i\).]}.\quad . . . . .
\par
Die 80 Gleichungen für die Inklination und den \(\AscendingNode\) ließen sich zu 40 für die Breite zusammenfassen. Ich glaube aber nicht, dass dadurch etwas gewonnen wird, denn wenn man die Elemente selbst den Störungen unterwirft, kann man ohne Bedenken dieselben gestörten Elemente als mehrere Monate lang gültig ansehen und braucht also nur einmal im Jahr für 6 Elemente die Störungen zu berechnen (insgesamt vielleicht etwa 300—400 Gleichungen). Bringt man dagegen bei den Elementen nur Säkularänderungen an und die periodischen Änderungen bei Breite, Länge und Radiusvektor (insgesamt vielleicht an die 200), so müssten diese innerhalb eines Jahres doch wohl für wenigstens 6 verschiedene Orte berechnet werden, um interpolieren zu können. Doch kann man dies in der \SourceLeaf{427}Folge machen, wie man will, wenn nur erst alle Störungen \textls[180]{in irgendeiner Form} vorliegen.
\par
